\section{Anharmonic and finite-temperature methods}\label{app:E.8}

The band gap of silicon near zero temperature is about \SI{60}{meV} smaller than the gap of the same crystal with its nuclei clamped at their equilibrium positions: the zero-point motion of the lattice closes it, and the measured values are 62 and \SI{64}{meV}~\cite[Table~VI]{ponce2015temperature}. Between zero and room temperature the gap closes by a further \SI{50}{meV} (\cref{sec:1.3}), and the phonons themselves shift and broaden. A detector at \SI{10}{mK} is deep in the zero-point regime, where all of these effects come from zero-point motion and none from thermal occupation. The anharmonic small parameter $\eta_{\mathrm{anh}}$ of \cref{sec:5.7} is small for silicon, and the lattice-side corrections are correspondingly small; the electron-side ones are not.

\subsection{Quasi-harmonic approximation}

In the quasi-harmonic approximation the free energy per cell, $\mathcal{F}$, at cell volume $\Omega$ is the harmonic one, with frequencies computed at that volume~\cite[Eq.~(4)]{baroni2010thermal},
\begin{equation}
  \mathcal{F}(\Omega,T) = E^{\mathrm{ad}}_0(\Omega) + \frac1N\sum_{\vect{Q}\nu}\frac{\omega_{\vect{Q}\nu}(\Omega)}{2} + \frac TN\sum_{\vect{Q}\nu}\ln\bigl[1 - \ee^{-\omega_{\vect{Q}\nu}(\Omega)/T}\bigr] ,
  \label{eq:E.8-qha}
\end{equation}
where $E^{\mathrm{ad}}_0(\Omega)$ is the clamped-nucleus energy per cell, the minimum of the energy surface of \cref{sec:5.1} at that volume, and the sums run over the $N$ points of the phonon grid. The equilibrium volume minimizes $\mathcal{F}$ at each $T$. Differentiating, the volume thermal expansion coefficient is
\begin{equation}
  \alpha_{\mathrm{th}} = \frac{1}{N B\Omega}\sum_{\vect{Q}\nu}\omega_{\vect{Q}\nu}\,\gamma^{\mathrm{G}}_{\vect{Q}\nu}\,\frac{\partial n_{\vect{Q}\nu}}{\partial T}, \qquad
  \gamma^{\mathrm{G}}_{\vect{Q}\nu} = -\frac{\Omega}{\omega_{\vect{Q}\nu}}\frac{\partial\omega_{\vect{Q}\nu}}{\partial\Omega} ,
  \label{eq:E.8-beta}
\end{equation}
with $B$ the isothermal bulk modulus and $\gamma^{\mathrm{G}}_{\vect{Q}\nu}$ the mode Grüneisen parameters~\cite[Eqs.~(6) and~(9)]{baroni2010thermal}. A perfectly harmonic crystal has $\gamma^{\mathrm{G}} = 0$ and no thermal expansion. The calculation is a set of harmonic calculations: DFPT phonons at five to seven volumes around the minimum, a fit of $E^{\mathrm{ad}}_0(\Omega)$ and of each $\omega_{\vect{Q}\nu}(\Omega)$, and a minimization of \cref{eq:E.8-qha}. The cost is that number of DFPT runs. In silicon the transverse acoustic branch has negative Grüneisen parameters, and its large weight in the density of states at low frequency makes the thermal expansion negative at low temperature, a behavior the quasi-harmonic approximation reproduces~\cite[Sect.~5.1]{baroni2010thermal}.

The approximation holds as long as the frequencies depend on the volume but not on the amplitude of vibration, that is, as long as the phonons are well-defined harmonic excitations at each volume. It fails for modes that are unstable in the harmonic approximation, near structural transitions and near melting, where the next two methods apply. The zero-point term of \cref{eq:E.8-qha} expands the lattice at $T = 0$, and this expansion changes the gap: for silicon the zero-point lattice expansion alone shifts the gap by $+\SI{9}{meV}$~\cite[Table~S2]{miglio2020predominance}, opposite to the electron--phonon renormalization. The same Grüneisen parameters also set how sensitive a phonon calculation is to the lattice constant it uses: by \cref{eq:E.8-beta}, a relative error $\delta a/a$ in the lattice constant shifts each frequency by $\delta\omega_{\vect{Q}\nu}/\omega_{\vect{Q}\nu} = -3\gamma^{\mathrm{G}}_{\vect{Q}\nu}\,\delta a/a$. In silicon the transverse acoustic modes, with $\gamma^{\mathrm{G}} < 0$, and the optical modes, with $\gamma^{\mathrm{G}} > 0$, move in opposite directions, so that no single choice between the theoretical and the measured lattice constant improves both. A calculation that needs the transverse acoustic branch accurately, as the phonon-assisted thresholds of \cref{chap:13} do, has to test both lattice constants and quote the spread.

\subsection{Self-consistent harmonic approximation}

When the anharmonicity is not small, the harmonic expansion about the minimum is the wrong starting point, and the self-consistent harmonic approximation (SCHA) replaces it by the best harmonic description of the actual distribution of nuclei. It minimizes the Gibbs--Bogoliubov bound on the free energy (\cref{sec:5.7})~\cite[Eq.~(2)]{monacelli2021stochastic} over the thermal states of harmonic Hamiltonians,
\begin{equation}
  H_{\mathrm{trial}} = (\text{nuclear kinetic energy}) + \frac12\sum_{ab}(X_a - \bar X_a)\,\Phi^{(2)}_{\mathrm{aux},ab}\,(X_b - \bar X_b) ,
  \label{eq:E.8-scha}
\end{equation}
with $X_a$ the $3N_{\mathrm{at}}$ Cartesian coordinates of the nuclei of a supercell of $N_{\mathrm{at}}$ atoms, and the centroids $\bar X_a$ and the auxiliary force constants $\Phi^{(2)}_{\mathrm{aux}}$ as variational parameters~\cite[Eqs.~(4) and~(7)]{monacelli2021stochastic}. Minimizing the bound over $\bar X$ and $\Phi^{(2)}_{\mathrm{aux}}$ gives the SCHA free energy, the equilibrium positions at temperature $T$, and an effective harmonic problem whose quantum zero-point motion is included by construction. At the minimum $\Phi^{(2)}_{\mathrm{aux}}$ satisfies the self-consistency condition of \cref{sec:5.7}: it equals the average of the true second derivatives of the energy surface over the Gaussian distribution it generates.

The averages are computed stochastically. A population of $N_{\mathrm{conf}}$ configurations is drawn in a supercell from the Gaussian distribution of the current trial Hamiltonian, each together with its mirror image about the centroids to reduce the statistical error; DFT supplies the energy, forces and stress of each; and the gradients of $\mathcal{F}$ with respect to $\bar X$ and $\Phi^{(2)}_{\mathrm{aux}}$ follow from these forces. As the parameters move, the same configurations are reweighted by the ratio of the new to the old Gaussian density, until the effective number of samples, $N_{\mathrm{eff}} = (\sum_jw_j)^2/\sum_jw_j^2$ for weights $w_j$, drops below about half of $N_{\mathrm{conf}}$; then a new population is drawn~\cite[Eqs.~(21), (38) and~(39)]{monacelli2021stochastic}. The cost is entirely in the DFT forces: a population holds several hundred to several thousand configurations, and most materials converge within three populations~\cite[Sect.~III\,B]{monacelli2021stochastic}. The convergence parameters are the supercell size, $N_{\mathrm{conf}}$, and the threshold on $N_{\mathrm{eff}}/N_{\mathrm{conf}}$.

The auxiliary frequencies of $\Phi^{(2)}_{\mathrm{aux}}$ are not the physical phonon energies. They describe the Gaussian distribution of the nuclei; the frequencies measured by a probe at small energy follow from the curvature of the free energy with respect to the centroids, which adds a term with third-order force constants~\cite[Sect.~IV]{monacelli2021stochastic}. Mistaking one for the other is the most common error in reading SCHA output. For silicon the SCHA changes the phonons little; it is needed for light atoms, soft modes and quantum crystals, the targets of \cref{sec:11.5}.

\subsection{Temperature-dependent effective potential}

The temperature-dependent effective potential (TDEP) method fits effective force constants to forces sampled at temperature $T$. Molecular dynamics in a supercell of $N_{\mathrm{at}}$ atoms supplies, at each of $N_t$ time steps, the displacements $\vect{u}_i$ and the DFT forces $\vect{F}_i$ on every atom $i$ of the supercell. The effective second-order force constants $\tilde\Phi^{(2)}_{ij}$ are those that minimize the mismatch
\begin{equation}
  \Delta_{\mathrm{fit}} = \frac{1}{N_t}\sum_{t,i}\Bigl|\vect{F}_{i}(t) + \sum_j\tilde\Phi^{(2)}_{ij}\,\vect{u}_j(t)\Bigr|^2 ,
  \label{eq:E.8-tdep}
\end{equation}
since a harmonic energy surface gives $\vect{F}_i = -\sum_j\tilde\Phi^{(2)}_{ij}\vect{u}_j$. This is a linear least-squares problem, overdetermined once $N_t > 3N_{\mathrm{at}}$~\cite[Eq.~(3), with the opposite sign convention for the force constants]{hellman2011lattice}. Symmetry reduces the unknowns to a few per coordination shell, and the same fit extends to third and fourth order. The fitted constants absorb the anharmonicity sampled at that temperature into effective harmonic ones; the phonons and free energy follow as in the harmonic case~\cite{hellman2013temperature}. The convergence parameters are the number of time steps, which Ref.~\cite[Fig.~1]{hellman2011lattice} tracks through the change in the fitted constants, the supercell size, and the range of the force constants kept in the fit.

Classical molecular dynamics samples a classical distribution, with mean squared displacement $T/(M_\kappa\omega_{\vect{Q}\nu}^2)$ per mode, while the quantum harmonic value is $(2n_{\vect{Q}\nu} + 1)/(2M_\kappa\omega_{\vect{Q}\nu})$. Their ratio is $2T/\omega_{\vect{Q}\nu}$ for $T \ll \omega_{\vect{Q}\nu}$. For silicon phonons of 10 to \SI{60}{meV} at \SI{10}{mK}, the ratio is below $2\times10^{-4}$, so a classical sample describes an almost static lattice and misses the zero-point motion that sets every renormalization at detector temperatures. At low temperature the configurations must be drawn from the quantum distribution instead, as the SCHA does, or from path-integral molecular dynamics (\cref{app:E.11}). The cost of TDEP is that of the molecular dynamics, $N_t$ force calculations of which only those separated by more than a correlation time are independent.

Third-order force constants also give the three-phonon linewidths $\gamma^{\mathrm{anh}}_{\vect{Q}\nu}$ of \cref{sec:5.7}. They are fitted in the same way, or computed by finite differences in supercells, one supercell calculation per pair of displaced atoms allowed by symmetry, with a cutoff radius on the interacting triplets. The linewidth then needs a sum over pairs of final phonons on a dense $\vect{Q}$ grid with an energy delta function, and the same rules for broadening and grid spacing apply as for electronic spectra (\cref{app:C.2}).

\subsection{Allen--Heine--Cardona calculations}

The AHC theory of \cref{sec:6.3} gives the shift of a band energy at second order in the electron--phonon coupling, summed over all phonons with the thermal weight $2n_{\vect{Q}\nu} + 1$. In the adiabatic form, which neglects the phonon energy next to electronic energy differences, the Fan--Migdal part is
\begin{equation}
  \Delta\varepsilon^{\mathrm{FM}}_{n\vect{k}}(T) = \frac{1}{N}\sum_{\vect{Q}\nu}(2n_{\vect{Q}\nu} + 1)\sum_m\frac{|g_{mn\nu}(\vect{k},\vect{Q})|^2}{\varepsilon_{n\vect{k}} - \varepsilon_{m,\vect{k}+\vect{Q}}} ,
  \label{eq:E.8-fan}
\end{equation}
with the vertex of \cref{eq:A.6-g}, and the Allen--Heine part adds the diagonal element of the second-order vertex $g^{\mathrm{AH}}$ with the same weight. In the rigid-ion approximation the second-order vertex is rewritten through first-order matrix elements at $\vect{Q} = 0$, using the translation invariance of the crystal, so that no second derivative of the potential is computed~\cite[Eqs.~(11) and~(15)]{ponce2015temperature} (\cref{sec:6.3}). The AHC literature calls this term the Debye--Waller term; in this book Debye--Waller refers only to $\ee^{-2W_\kappa}$ (\cref{app:A.6}). The nonadiabatic form keeps $\pm\omega_{\vect{Q}\nu}$ in the denominators of \cref{eq:E.8-fan}, separately for phonon emission and absorption, and gives the renormalization a lifetime as well~\cite[Eq.~(16)]{ponce2015temperature}.

The sum over $m$ runs over all bands. In practice it is split: the bands within an active window are summed explicitly, and the contribution of all higher bands is obtained from a Sternheimer equation for the first-order change of the wavefunction projected on the space above the window, as DFPT does for phonons (\cref{app:E.4}). The split is exact in the adiabatic limit; in the nonadiabatic form the Sternheimer part drops the phonon energy, and the active window must be large enough for that to be harmless~\cite[below Eq.~(16)]{ponce2015temperature}.

The sum over $\vect{Q}$ is the hard part. The Fan--Migdal denominators vanish when $\varepsilon_{m,\vect{k}+\vect{Q}} = \varepsilon_{n\vect{k}}$, which happens for $m = n$ at $\vect{Q} \to 0$ and wherever another band is degenerate with the state. Near these points the integrand is singular, and the result converges slowly with the number $N_Q^3$ of grid points, at a rate that depends on the band and on the form of the theory; a small imaginary part $\ii\delta$ in the denominators smooths the integrand but introduces a bias that must be extrapolated to $\delta \to 0$ together with $1/N_Q \to 0$~\cite[Sect.~III and Table~II]{ponce2015temperature}; the broadening is a numerical device, never quoted as a lifetime. A nonzero Born effective charge makes the vertex diverge as $1/Q$ at small $Q$. In a nonpolar crystal the charges vanish by symmetry, but a DFPT calculation always leaves small residual charges that violate the charge-neutrality sum rule, and these must be removed before the $\vect{Q}$ sum, or the integral diverges at band extrema~\cite[Sect.~III]{ponce2015temperature}. In a polar crystal the divergence is physical: the adiabatic integral diverges, and only the nonadiabatic form gives a finite result~\cite{ponce2015temperature}. For polar materials the nonadiabatic, Fröhlich-like part often dominates the zero-point renormalization~\cite{miglio2020predominance}.

The cost is that of DFPT at every $\vect{Q}$ of a dense grid: for silicon, Ref.~\cite{ponce2015temperature} goes up to a $100\times100\times100$ grid, 22\,776 irreducible points, to converge the values of its Table~VI. Wannier interpolation of the vertices (\cref{app:E.5}) reduces this to DFPT on a coarse grid. The results for silicon are in \cref{we:E-zpr}.

\subsection{Special displacements}

The perturbative route needs vertices and a converged sum over phonon wavevectors. The non-perturbative alternative computes the electronic property directly in a supercell whose nuclei are displaced, and averages over displacements. In the adiabatic and harmonic approximations, the Williams--Lax expression for a temperature-dependent property, here the imaginary part of the dielectric function, is a Gaussian average over the normal coordinates $x_\nu$ of the supercell,
\begin{equation}
  \varepsilon_2(\omega;T) = \prod_{\vect{Q}\nu}\int\dd x_{\vect{Q}\nu}\,\frac{\ee^{-x_{\vect{Q}\nu}^2/2\sigma_{\vect{Q}\nu}^2}}{\sqrt{2\pi}\,\sigma_{\vect{Q}\nu}}\;\varepsilon_2(\omega;\{x\}), \qquad \sigma_{\vect{Q}\nu}^2 = \frac{2n_{\vect{Q}\nu} + 1}{2\omega_{\vect{Q}\nu}} ,
  \label{eq:E.8-wl}
\end{equation}
where the $x_{\vect{Q}\nu}$ are the mass-weighted normal coordinates of the supercell, labeled by the wavevectors $\vect{Q}$ of the primitive-cell phonons it contains, and $\varepsilon_2(\omega;\{x\})$ is computed with the nuclei frozen at the displaced positions~\cite[Eqs.~(3) and~(4), with mass-weighted coordinates]{zacharias2016one}. The same formula holds for a band energy or any other property computed at fixed nuclei.

The average needs only one deterministic configuration. In a large supercell, the average of a smooth function over the Gaussian distribution of \cref{eq:E.8-wl} is dominated by configurations in which every normal coordinate has the magnitude $\sigma_{\vect{Q}\nu}$, and the cross terms between pairs of modes cancel if the signs alternate in a suitable order~\cite[Sect.~III]{zacharias2016one}. The resulting special, or ZG, displacement of nucleus $\kappa$ in cell $\vect{R}$ of a supercell of $N$ primitive cells is
\begin{equation}
  \vect{u}^{\mathrm{ZG}}_{\vect{R}\kappa} = \sum_{\vect{Q}\in\mathcal{B},\nu}\varsigma_{\vect{Q}\nu}\left[\frac{2n_{\vect{Q}\nu} + 1}{2NM_\kappa\omega_{\vect{Q}\nu}}\right]^{1/2}2\,\Real\bigl[\ee^{\ii\vect{Q}\cdot\vect{R}}\,\vect{e}_{\kappa\nu}(\vect{Q})\bigr] ,
  \label{eq:E.8-zg}
\end{equation}
in the conventions of \cref{eq:A.6-u}~\cite[Eq.~(2)]{zacharias2020theory}. The set $\mathcal{B}$ contains one wavevector from each pair $\pm\vect{Q}$ of the grid commensurate with the supercell, and the signs $\varsigma_{\vect{Q}\nu} = \pm1$ alternate from one mode to the next in a prescribed order. Wavevectors equal to their own time-reversal partners, such as $\vect{Q} = 0$, are left out: they describe standing waves that break the crystal symmetry, and their weight vanishes as the supercell grows~\cite[Sects.~II and~VII\,B]{zacharias2020theory}. The mean squared displacement of each atom, averaged over the supercell, then equals its thermal value from \cref{eq:A.6-u}. The method needs the phonons of the primitive cell on the commensurate grid, from DFPT, and one DFT calculation of the displaced supercell with an electronic $\vect{k}$ sampling appropriate to its size.

The convergence parameter is the supercell size. Including or excluding the wavevectors with $\vect{Q} = -\vect{Q}$ modulo $\vect{G}$ brackets the converged value from above and below, and the width of the bracket is a direct error estimate~\cite[Sect.~VII\,B]{zacharias2020theory}. With the $\vect{Q} = -\vect{Q}$ modes included, degenerate bands such as the threefold valence maximum of silicon split, and the splitting is an artifact~\cite[Sect.~VII\,B]{zacharias2020theory}. The phases of the eigenvectors must be fixed consistently from mode to mode before the signs are assigned; Ref.~\cite[Sect.~II]{zacharias2016one} chooses the first nonzero component of each eigenvector positive. Since the displaced supercell is a static calculation, it needs no sum over phonon wavevectors with vanishing denominators, and the adiabatic divergence of polar materials is absent: for GaAs the zero-point renormalization converges smoothly with supercell size, though a $6\times6\times6$ supercell, 432 atoms, is needed to come within ten percent~\cite[Sect.~VII]{zacharias2020theory}. The cost is one supercell calculation, which scales as the cube of the number of atoms for the diagonalization; an $8\times8\times8$ supercell of silicon holds 1024 atoms and costs about $8^3 = 512$ times more than a $4\times4\times4$ one at fixed accuracy per atom.

\paragraph{Phonon-assisted absorption.} With \cref{eq:E.8-wl} applied to $\varepsilon_2$, a single ZG supercell gives the temperature-dependent absorption spectrum including indirect transitions. The displaced supercell breaks translational symmetry, so that a photon can connect states that differ in crystal momentum by a phonon wavevector of the supercell grid. For silicon at room temperature, an $8\times8\times8$ supercell with 40 random $\vect{k}$ points, a scissor of \SI{0.74}{eV} and a Gaussian broadening of \SI{30}{meV} reproduces the measured absorption coefficient over seven orders of magnitude~\cite[Sect.~II and~VII]{zacharias2016one}. Being adiabatic, the method does not distinguish the absorption of a phonon from its emission; it cannot reproduce the two thresholds $\Eg \pm \omega_{\vect{Q}\nu}$ of the Hall--Bardeen--Blatt theory near the onset~\cite[Sect.~II]{zacharias2016one}, the region that sets the phonon-assisted thresholds of \cref{chap:13}. At \SI{10}{mK} the phonon-absorption branch is absent, and the onset is at $\Eg + \omega_{\vect{Q}\nu}$ (\cref{sec:6.5}). Special displacements therefore give the spectrum and its renormalization away from the edge, while the edge itself needs the nonadiabatic treatment with explicit phonon energies of \cref{app:E.5} (\cref{sec:13.7}).

\paragraph{Zero-point renormalization from special displacements.} The band energies of the displaced supercell, unfolded onto the primitive zone as a spectral function, give the renormalized bands $\varepsilon^{\mathrm{ren}}_{n\vect{k}}(T)$ directly. Unlike the perturbative sum, this includes all orders in the displacements, within the adiabatic approximation, and needs no electron--phonon vertices. It misses the nonadiabatic correction; in silicon the adiabatic and nonadiabatic AHC renormalizations differ by \SI{8}{meV}~\cite[Table~VI]{ponce2015temperature}.

\begin{workedexample}{The zero-point renormalization of the silicon gap}\label{we:E-zpr}
For the indirect gap of silicon at $T = 0$, perturbative AHC theory with the local-density approximation, Sternheimer equations for the high bands and grids up to $100\times100\times100$ phonon wavevectors, extrapolated to $\delta \to 0$, gives $-\SI{64.3}{meV}$ in the adiabatic form and $-\SI{56.2}{meV}$ in the nonadiabatic one~\cite[Table~VI]{ponce2015temperature}. The ZG displacement in an $8\times8\times8$ supercell, 1024 atoms, gives $-\SI{57}{meV}$, made of a rise of the valence maximum by \SI{35}{meV} and a drop of the conduction minimum, near $0.87$ of the way from $\mathrm{\Gamma}$ to $\mathrm{X}$, by \SI{22}{meV}; a $10\times10\times10$ supercell brackets the converged value between 55 and \SI{65}{meV}~\cite[Sect.~VII\,B]{zacharias2020theory}. A $4\times4\times4$ supercell, 128 atoms, is already within ten percent of the converged value~\cite[Sect.~VII]{zacharias2020theory}. The measured renormalization is $-62$ or $-\SI{64}{meV}$, depending on how the data are extrapolated~\cite[Table~VI]{ponce2015temperature}.

The measured value contains both the electron--phonon renormalization and the lattice expansion, which contributes $+\SI{9}{meV}$ (\cref{eq:E.8-qha}). The lattice term is therefore added to the computed fixed-volume value before the two are compared~\cite{miglio2020predominance}: $-56 + 9 = -\SI{47}{meV}$ for the nonadiabatic AHC result and $-57 + 9 = -\SI{48}{meV}$ for the ZG displacement, against $-62$ to $-\SI{64}{meV}$. Equivalently, the electron--phonon part implied by the measurement is $-71$ to $-\SI{73}{meV}$, and the calculations underestimate it by about \SI{20}{\percent}, a shortfall that \cref{sec:6.3} traces to the overscreening of the vertex by the small Kohn--Sham gap. The cost differs by orders of magnitude: tens of thousands of DFPT calculations in the primitive cell for the perturbative route, against DFPT phonons on an $8\times8\times8$ grid and one 1024-atom supercell for the special displacement. A measured low-temperature gap, such as the \SI{1.17}{eV} of silicon (\cref{sec:1.3}), already contains the zero-point renormalization (\cref{sec:3.4}), and a computed renormalization is added only to a clamped-nucleus gap, from $GW$ or from a hybrid functional. Between \SI{10}{mK} and \SI{1}{K} all occupations $n_{\vect{Q}\nu}$ of silicon are of order $\ee^{-\omega_{\vect{Q}\nu}/T}$ with $\omega_{\vect{Q}\nu}/T \gtrsim 10^2$ for every phonon above \SI{10}{meV}, and the thermal part of the renormalization is negligible; only the acoustic modes at small $Q$ are thermally occupied, and their contribution to \cref{eq:E.8-fan} vanishes with their phase space. At those temperatures the populations of the final states still depend on temperature, and their energies do not.
\end{workedexample}
